\documentclass[12pt, a4paper]{article}

\usepackage[T2A]{fontenc}
\usepackage[utf8]{inputenc}
\usepackage[bulgarian]{babel}
\usepackage{geometry}
\geometry{a4paper, margin=1in}
\usepackage{tabularx}
\usepackage{xltabular}
\usepackage{ragged2e}
\usepackage{booktabs}
\usepackage{hyperref}

\title{Задание 3: Тестване на прототип}
\author{}
\date{}

\begin{document}

\maketitle

\noindent\textbf{Дисциплина:} Проектиране на човеко-машинен интерфейс 2025-2026

\bigskip

\noindent\textbf{Участници в проекта:}
\smallskip

\noindent
\begin{tabularx}{\textwidth}{|l|X|l|X|l|}
\hline
\textbf{№} & \textbf{Име и фамилия} & \textbf{Факултетен №} & \textbf{Специалност} & \textbf{Курс} \\ \hline
1 & Виктор Койчев & 3MI0600164 & Софтуерно инженерство & 3-ти \\ \hline
2 & Емилиан Янев & 0MI0600273 & Софтуерно инженерство & 3-ти \\ \hline
3 & Жан Георгиев & 8MI0600398 & Софтуерно инженерство & 3-ти \\ \hline
4 & Даниел Николов & 0MI0600410 & Софтуерно инженерство & 3-ти \\ \hline
\end{tabularx}

\bigskip

\noindent\textbf{Име на група:}\\
HCI\_2026\_group\_3MI0600164\_0MI0600273\_8MI0600398\_0MI0600410

\bigskip

\noindent\textbf{Име на проекта:}\\
Система за управление и мониторинг на роботизиран склад за фармацевтични продукти (PharmaBot Logistics)

\bigskip
\rule{\textwidth}{0.4pt}

\section*{1. Методология на тестването}

Тестването на първоначалния (Low-Fidelity) HTML прототип бе проведено чрез метода на експертната оценка (Heuristic Evaluation), базирана на 10-те хевристики за ползваемост на Якоб Нилсен, както и чрез когнитивно обхождане (Cognitive Walkthrough) на основните потребителски сценарии от перспективата на Складов оператор и Шофьор на камион.

\section*{2. Оценка на интерфейса по хевристиките на Нилсен}

\textbf{1. Видимост на състоянието на системата (Visibility of system status)}\\
\textit{Реализация:} В десктоп дашборда е интегрирана 2D карта, която показва позицията на роботите в реално време. Състоянието им е ясно индикирано чрез мигащи анимации (в движение) и статични икони (зареждане).

\textbf{2. Връзка между системата и реалния свят (Match between system and the real world)}\\
\textit{Реализация:} Използвани са термини, познати на логистичните работници: „Партида“, „Сектор А“, „Товарна рампа“, „Студена верига“.

\textbf{3. Потребителски контрол и свобода (User control and freedom)}\\
\textit{Реализация:} При създаване на „Зона за опасност“ по погрешка, операторът може лесно да я премахне само с един клик. Модалните прозорци имат ясни опции за отказ/игнориране.

\textbf{4. Последователност и стандарти (Consistency and standards)}\\
\textit{Реализация:} Спазени са глобалните цветови стандарти: Зелено (успех), Жълто (внимание) и Червено (опасност).

\textbf{5. Превенция на грешки (Error prevention)}\\
\textit{Реализация:} В мобилното приложение за шофьори, бутоните за следващи стъпки са деактивирани и визуално заключени (overlay lock), докато не бъде завършена предходната стъпка.

\textbf{6. Разпознаване вместо запаметяване (Recognition rather than recall)}\\
\textit{Реализация:} При назначаване на нова задача, операторът не трябва да въвежда ръчно ID на робот – системата предлага падащо меню с наличните роботи и текущия им заряд.

\textbf{7. Гъвкавост и ефективност на използване (Flexibility and efficiency of use)}\\
\textit{Реализация:} За напреднали потребители е добавен бърз бутон за превключване в режим „Очертай зона“, който променя курсора и позволява светкавична реакция при разлив.

\textbf{8. Естетичен и минималистичен дизайн (Aesthetic and minimalist design)}\\
\textit{Реализация:} Мобилният интерфейс е лишен от излишна информация. Десктоп изгледът фокусира 80\% от екрана върху най-важното – картата.

\textbf{9. Разпознаване, диагностика и възстановяване от грешки (Help users recognize, diagnose, and recover from errors)}\\
\textit{Реализация:} При температурна аларма, системата изписва точния проблем: „Текуща температура 12.5°C (Макс: 8.0°C)“ и предлага директен екшън бутон: „Пренасочи към Карантина“.

\textbf{10. Помощ и документация (Help and documentation)}\\
\textit{Реализация:} Мобилният процес за предаване на пратката е разделен на изрични стъпки („Стъпка 1: Идентификация“, „Стъпка 2: Доказателство“), което служи като вградено ръководство за потребителя.

\section*{3. Открити проблеми при първото тестване и тяхното отстраняване (Итеративен дизайн)}

По време на тестването и когнитивното обхождане идентифицирахме 3 критични UX проблема. За да приложим принципите на итеративния дизайн, \textbf{веднага нанесохме корекции в HTML прототипа} (версия 0.2):

\begin{enumerate}
\item \textbf{Проблем:} \textit{Липса на потвърждение при изход (Нарушаване на хевристика 5 -- Превенция на грешки).} В първата версия линкът „Изход от системата“ не изискваше потвърждение, което създаваше риск от инцидентно отписване по време на критична операция.
\\
\textbf{Корекция в прототипа:} В \texttt{dashboard.html} добавихме модален прозорец \texttt{logoutModal}, който изисква изрично потвърждение, а в мобилната версия (\texttt{mobile-handoff.html}) интегрирахме Native Browser Confirm диалог, предупреждаващ, че „Текущият напредък по пратката ще бъде запазен“.

\item \textbf{Проблем:} \textit{Претрупване на 2D картата (Нарушаване на хевристика 8).} При наличие на множество роботи (движещи се, зареждащи се, сгрешени), картата за складовия оператор ставаше прекалено шумна визуално.
\\
\textbf{Корекция в прототипа:} В \texttt{dashboard.html} внедрихме интерактивен панел „Филтри на картата“. Чрез чекбоксове операторът вече може да скрива роботи, които се зареждат, за да фокусира вниманието си само върху движещите се машини или тези с проблем.

\item \textbf{Проблем:} \textit{Ниска видимост на открито за мобилното приложение (Accessibility / Контраст).} Шофьорите работят навън на товарните рампи. Първоначалният светлосив дизайн на неактивните бутони се оказа труден за разчитане на пряка слънчева светлина.
\\
\textbf{Корекция в прототипа:} В \texttt{mobile-handoff.html} направихме базовия контраст на рамките много по-осезаем. Допълнително добавихме бутон \textbf{"Контраст"}, който превключва приложението в специализиран тъмен/висококонтрастен режим (Dark Mode), изчиствайки отблясъците на екрана. Също така добавихме ясен визуален „катинар“ (overlayLock) с надпис „Изисква Стъпка 1“, за да е недвусмислено защо бутонът за камерата не работи.
\end{enumerate}

\end{document}